Run the rational DCTFTR phase first. Recursion soundness gives a unique finite trace. If it returns a successful rational identification, its expression belongs to \(\mathscr R(\mathcal P)\) and equals the causal query throughout \(\mathcal C(\mathcal P)\); relabelling this terminal as \(\operatorname{ID\_RAT}\) does not change its derivation. A certified uniform-zero terminal is likewise sound. Every other local terminal invokes the canonical fallback specified by the hybrid grammar.

Set
\[
S=\Theta_G^{\mathrm{can}}(\mathcal P).\tag{1}
\]
By the exact global canonical theorem and the assumption
\(\mathcal C(\mathcal P)\ne\varnothing\), \(S\) is a nonempty compact semialgebraic set whose target-cell range is exactly the broad SCM range on the original alphabets. All defining coefficients lie in the effectively presented real-algebraic field generated by the supplied menu values.

For a polynomial \(p\) on \(S\), its image is described by
\[
\varphi_p(t)\ \Longleftrightarrow\ \exists x\,[x\in S\ \wedge\ t-p(x)=0].\tag{2}
\]
For polynomials \(p,q\) with \(q>0\) on \(S\), the ratio image is described by
\[
\varphi_{p/q}(t)\ \Longleftrightarrow\ \exists x\,[x\in S\ \wedge\ q(x)>0\ \wedge\ q(x)t-p(x)=0].\tag{3}
\]
Effective real-algebraic quantifier elimination terminates on (2)--(3), returns equivalent quantifier-free formulas, and supplies a finite sign-invariant cylindrical decomposition with algebraic sample points. Since \(S\) is compact and the relevant functions are continuous, their images have attained endpoints. Ordering the one-dimensional output cells locates those endpoints exactly, and lifting their closed endpoint cells through (2) or (3) returns feasible attaining parameter samples. Projection polynomials, isolating intervals, sign tables, truth propagation, and direct substitution give finite exact certificates. The canonical ordering in the hybrid grammar makes the proof-producing run unique; correctness does not depend on uniqueness of a possible CAD representation.

For an unconditional target use
\(p=\Phi_{T,\mathrm{can}}^\star\) in (2). Exact range equality implies that equal endpoints occur exactly when the target is identified, producing \(\operatorname{ID\_CAN}\); unequal endpoints lift to two original-alphabet same-menu SCMs and produce \(\operatorname{PAIR\_CAN}\).

For a conditional query, first apply (2) to
\(p=\Phi_{E,\mathrm{can}}^\star\). Nonnegativity of evidence makes the following cases exhaustive:
\[
\begin{array}{ll}
\max_S p=0,&\text{evidence is zero in every compatible model};\\
\min_S p=0<\max_S p,&\text{there are respectively zero- and positive-evidence witnesses};\\
\min_S p>0,&\text{evidence is uniformly positive}.
\end{array}\tag{4}
\]
The first two cases return \(\operatorname{ZERO\_EVIDENCE}\) and
\(\operatorname{NOT\_UNIFORMLY\_DEFINED}\), respectively, before any division. Only in the third case apply (3) with
\[
p=\Phi_{B_\rho\cap E,\mathrm{can}}^\star,\qquad
q=\Phi_{E,\mathrm{can}}^\star.\tag{5}
\]
Equal sharp ratio endpoints give \(\operatorname{ID\_CAN}\); unequal endpoints give \(\operatorname{PAIR\_CAN}\), with the two lifted SCMs attaining unequal conditional values. Exact global range equality proves the stated exactness clauses in every case.

The rational phase terminates, and whenever it does not settle the query, the finite quantifier-elimination phase terminates. Thus the combined run is finite and returns exactly one of the grammar's exhaustive terminals. Every witness stays on \(\mathcal X\), agrees with all supplied laws, and is checked by exact real-algebraic arithmetic. Hence every unresolved local trace is completed by identification, an undefinedness certificate, or an explicit paired-model certificate; local row-face variation alone is never used as broad-model nonidentification.